\centerline{\bf \Huge Показатели}
\centerline{\bf \Huge Решения}

\problem{1} Докажите, что если $p$ --- простое, то у числа $2^p-1$ все
простые делители больше $p$.

\textbf{Решение.} Предположим противное: пусть $2^p\equiv_q 1$ и $q
\leqslant p$ --- некоторый простой делитель числа $2^p-1$. Пусть
также $T$ --- порядок числа 2 по модулю $q$. Тогда $2^T\equiv_q 1$ и
$T\mid p$, откуда либо $T=1$, либо $T=p$. В первом случае получаем,
что $2\equiv_p 1$, что невозможно, а во втором случае
получаем, что $T = p\mid q-1$ по малой теореме Ферма, откуда
$p\leqslant q-1$ --- противоречие.

\problem{2} Докажите, что если $p$ --- простое, то любой простой делитель
числа $2^p-1$ имеет вид $2kp+1$ для некоторого натурального $k$.

\textbf{Решение.} Пусть $q$ --- простой делитель числа $2^p-1$. Если $T$ ---
порядок числа 2 по модулю $q$, то $T\mid p$, откуда $T=1$ или $T=2$.
В первом случае получаем, что $2\equiv_q 1$, что невозможно,
а во втором --- что $p\mid q-1$ по малой теореме Ферма. Значит,
$q-1=np$. Ясно, что $q$ нечетно, поэтому $n=2k$ четно и $q=2kp+1$.

\problem{3} Докажите, что все простые делители чисел Ферма $2^{2^n}+1$ имеют вид
$2^{n+1}\cdot x+1$. (Именно так Эйлер нашел простой делитель у числа
$2^{2^5}+1$.)

\textbf{Решение.} Пусть $2^{2^n}\equiv_p -1$. Тогда $2^{2^{n+1}}\equiv_p
1$. Пусть $T$ --- порядок числа $2$ по модулю $p$.
Получаем, что $T\mid 2^{n+1}$, т.е. $T=2^k$ при некотором
$k\in\mathbb{N}$. Если $k\leqslant n$, то $1\equiv_p
(2^{2^k})^{2^{n-k}}=2^{2^n}\equiv_p -1$ и $p=2$, что
невозможно. Значит, $T=2^{n+1}$ и по малой теореме Ферма
$2^{n+1}=T\mid p-1$, что и требовалось доказать.

\problem{4} Докажите, что для любого натурального $n$ число $2^n-1$ не
делится на $n$.

\textbf{Решение.} Предположим, что для некоторого натурального $n$ число $2^n-1$
делится на $n$. Рассмотрим наименьший простой делитель $p$ числа $n$
и порядок $T$ числа 2 по модулю $p$. Тогда $2^n\equiv_p 1$,
откуда $T\mid n$ и $T\mid p-1$. Но тогда $T < p$ --- делитель $n$,
меньший $p$. Значит, $T=1$ и $2\equiv_p 1$ --- противоречие.

\problem{5} Существует ли набор натуральных чисел $(n_1,\ldots,n_k)$,
такой, что $n_i\mid 2^{n_{i+1}}-1$ для всех $i=1$, \ldots, $k$ (мы
считаем, что $n_{k+1}=n_1$)?

\textbf{Решение.} Пусть $p$ --- наименьший простой делитель среди всех делителей
чисел $n_1$, \ldots, $n_k$. Тогда $p\mid n_i\mid 2^{n_{i+1}}-1$.
Значит, $2^{n_{i+1}}\equiv_p 1$. Пусть $T$ --- порядок числа
$2$ по модулю $p$. Тогда $T < p$ и $T\mid n_{i+1}$. Если $T>1$, то у
числа $n_{i+1}$ существует простой делитель (входящий в разложение
$T$), меньший $p$, что невозможно. Значит, $T=1$ и $2\equiv_p 1$, что невозможно. Значит, искомых наборов не существует.

\problem{6} Найдите все тройки простых чисел $(p,q,r)$, такие,
что $p \mid q^r+1$, $q \mid r^p+1$, $r \mid p^q+1$.

\textbf{Решение.} Если $q^r\equiv_p -1$, то $q^{2r}\equiv_p 1$. Если $T$ --- порядок числа $q$ по модулю $p$, то $T\mid 2r$, откуда $T\in \{1,2,r,2r\}$. 

Если $T=1$, то $q\equiv_p 1$, $1\equiv_p q^r\equiv_p -1$ и $p=2$. 

Если $T=r$, то опять $1\equiv_p q^r\equiv_p  -1$ и $p=2$. 

Если $T=2$, то $q^2\equiv_p 1$, откуда $p\mid q^2-1$.

Если $T=2r$, то по малой теореме Ферма $2r=T\mid p-1$. 

Аналогичные рассуждения можно провести и для остальных делимостей.

Ясно, что все числа $p$, $q$, $r$ различны. Предположим, что все
числа $p$, $q$, $r$ больше 2. Тогда либо $2r\mid p-1$,
либо $p\mid q^2-1$. В первом случае получаем, что $0\equiv_r
p^q+1\equiv_r 1+1=2$, т.е. $r=2$ --- противоречие. Во втором
случае либо $p\mid q-1$, либо $p\mid q+1$. Если $p\mid q-1$, то
$0\equiv_p q^r+1\equiv_p 1+1=2$ и $p=2$ --- противоречие.
Значит, $p\mid q+1$. Аналогично, $q\mid r+1$ и $r\mid p+1$. Но тогда
$p\leqslant q\leqslant r\leqslant p$ и $p=q=r$ --- противоречие.

Итак, без ограничения общности $p=2$. Тогда $q$ и $r$ нечетны. Тогда либо $2q\mid r-1$, либо $r\mid
2^2-1=3$. В первом случае $r\equiv_q 1$ и $0\equiv_q
r^2+1\equiv_q 1+1=2$, т.е. $q=2$, что невозможно. Значит,
$r=3$ и $q=5$.

\claim{Ответ} $(2;5;3)$, $(3;2;5)$, $(5;3;2)$.
\problem{7} Найти все пары простых чисел $(p;q)$, таких, что $pq \mid
(5^p-2^p)(5^q-2^q)$.

\textbf{Решение.} Пусть $p\mid 5^p-2^p$. Тогда $0\equiv_p 5^p-2^2\equiv_p 5-2=3$ по малой теореме Ферма, откуда $p=3$. Значит, $q\mid (5^3-2^3)(5^q-2^q)$ и либо $q=3$, либо $q\mid 5^3-2^3=117$ и $q=13$. Аналогично, если $q=3$, то $p=3$ или $p=17$. Теперь предположим, что $p\neq 3$ и $q\neq 3$. Тогда $p\mid 5^q-2^q$. Перепишем это в виде $5^q\equiv_p 2^q$ и разделим это сравнение на $2^q$. Это можно сделать, т.е. $2^q$ --- делитель единицы по модулю $p$. Пусть $a=5/2\in \mathbb{Z}_p$, тогда $a^q\equiv_p 1$. Обозначим через $T$ порядок числа $a$ по модулю $p$. Тогда или $T=1$, или $T=q$. При $T=1$ получаем, что $a\equiv_p 1$ и $5\equiv_p 2$, т.е. $p=3$ --- противоречие. При $T=q$ из малой теоремы Ферма получаем, что $q\mid p-1$. Аналогично, получаем, что $p\mid q-1$. Но тогда $q < p$ и $p < q$ --- противоречие.

\claim{Ответ} $(3;3)$, $(3;13)$, $(13;3)$.

\problem{8} Найти все пары простых чисел $p$, $q$, таких, что $pq \mid
2^p+2^q$.

\textbf{Решение.} Заметим, что если $p=2$, то $q\mid 2 + 2^{q-1}$ и либо $q=2$,
либо $q=3$. Аналогично, при $q=2$ получаем, что $p=2$ или 3. Будем
считать, что $p$ и $q$ не равны 2. Пусть $p-1=2^k\cdot n$ и
$q=2^l\cdot m$, где $n$, $m$ нечетны. Без ограничения общности будем
считать, что $k\geqslant l$. Положим $x=2^n$, тогда согласно малой
теореме Ферма, $2^{p-1}\equiv_q -1$. Значит, порядок числа
$x$ по модулю $q$ равен $2^{k+1}$, поскольку $x^{2^k}=2^{p-1}\equiv_q
-1$ и $x^{2^{k+1}}=2^{2(p-1)}\equiv_q 1$. Отсюда
следует, что $2^{k+1}\mid q-1 = 2^l\cdot m$, т.е. $k+1\leqslant
l\leqslant k$ --- противоречие.

\problem{9} Найдите все натуральные $n$, для которых числа
$n$ и $2^n + 1$ имеют один и тот же набор простых делителей.

\textbf{Решение.} Заметим, что $n$ нечетно, т.к. в противном случае $2^n+1$ делится на 2, что неверно. Теперь рассмотрим наибольший простой делитель $p$ числа $n$. Докажем, что если $p>3$, то у числа
$2^p+1$ все простые делители больше $p$. В самом деле, пусть
$2^p\equiv_q -1$, где $q$ простое. Тогда $2^{2p}\equiv_q 1$, поэтому показатель числа $2$ по модулю $q$ равен 2 или $2p$. В первом случае получаем, что $q=3$, а во втором --- что $q> 2p>p$.

Рассмотрим наибольший простой делитель $p$ числа $n$. Тогда $p$ --- наибольший простой делитель числа $2^n+1$. С другой стороны, т.к. $n$ нечетно, что $2^p+1\mid 2^n+1$, а у числа $2^p+1$, как мы доказали, все простые делители больше $p$, если $p>3$. 

Значит, $p=3$ и $n=3^k$. Тогда $2^{n}+1 = 3^m$. Ясно, что $m$ четно и равно $2\ell$ (сравним обе части по модулю $4$). Тогда $3^m-1=(3^\ell-1)(3^\ell+1)=2^n$. Получается, что числа $3^\ell-1$ и $3^\ell+1$ являются степенями двойки. Но между ними разница 2, поэтому они равны 2 и 4 соответственно. Окончательно находим отсюда $\ell=1$, $m=2$ и $n=3$.

\claim{Ответ} 3.

